\documentclass[11pt]{article}
\usepackage{vinotheca}
\usepackage{pdflscape}

\begin{document}

\vinothecaTitle{TasteRank Explorer}{Data Appendix}{Source Data, Tables, and Reference Materials}
\vinothecaAuthor{Jure Skarabot}{May 2026}
\vinothecaCompanions{Summary \textperiodcentered{} Technical Appendix \textperiodcentered{} Methods Primer \textperiodcentered{} Data Appendix \textperiodcentered{} Grape Reference}

\section*{Table of Contents}
\begin{itemize}[leftmargin=2em,itemsep=0.1em,topsep=0.1em]
\item D.1 Sensory Dimension Key
\item D.2 Community Summary
\item D.3 Complete Variety Data (Ranks 1--101)
\end{itemize}

\section*{D.1 Sensory Dimension Key}
The following abbreviations are used in the data tables. All scores are integers on a 0--5 scale representing canonical varietal typicity across thirteen sensory dimensions drawn from established wine-tasting methodology (UC Davis, Peynaud, and related structured tasting frameworks).

\begin{center}
\begin{tabular}{@{}lll@{}}
\toprule
Abbr.\ & Dimension & Description \\
\midrule
Col & Color Depth         & Intensity of pigmentation (pale to opaque) \\
Aro & Aromatic Intensity  & Strength of nose from low to pronounced \\
Flo & Floral              & Degree of floral aromatic contribution \\
Frt & Fruit Ripeness      & Ripe / dried fruit vs.\ underripe / green fruit \\
Hrb & Herbal / Earthy     & Green herb, earth, mineral notes \\
Spk & Spice / Oak         & Oak-derived and spice character \\
Acd & Acidity             & Perceived acidity from low to high \\
Tan & Tannin              & Tannin level (reds); 0 for most whites \\
Bod & Body                & Weight and texture on palate \\
Alc & Alcohol             & Perceived alcohol warmth \\
Fla & Flavor Intensity    & Palate flavour concentration \\
Fin & Finish              & Length and persistence of aftertaste \\
Cpx & Complexity          & Layeredness and evolution of character \\
\bottomrule
\end{tabular}
\end{center}

\section*{D.2 Community Summary}
Six communities were detected by the Clauset--Newman--Moore greedy modularity optimisation algorithm. The table below summarises each community's character, size, colour composition, and hub varieties.

\begin{center}
\begin{tabular}{@{}clrlp{6.5cm}@{}}
\toprule
ID & Character & Size & R / W & Hub Varieties \\
\midrule
C0 & Full-bodied Mediterranean Reds   & 31 & 31 / 0 & Sagrantino, Nero d'Avola, Lagrein, Negroamaro, Plavac Mali \\
C1 & Light / Aromatic Cross-boundary  & 21 & 4 / 17 & Assyrtiko, Falanghina, Albari\~no, Riesling, Kadarka \\
C2 & Mid-weight Structured Reds       & 18 & 18 / 0 & Limniona, St.\ Laurent, Zweigelt, Blaufr\"ankisch, Sangiovese \\
C3 & Rich Textural Whites             & 13 & 0 / 13 & Marsanne, Fiano, Godello, Pinot Gris, Chardonnay \\
C4 & Mineral / Crisp Whites           & 10 & 0 / 10 & Greco, Friulano, Verdicchio, Verdejo, Vermentino \\
C5 & Lean Neutral \& Aromatic Whites  &  8 & 0 / 8  & Gew\"urztraminer, Petit Manseng, Muscat Blanc, Malvasia \\
\bottomrule
\end{tabular}
\end{center}

\begin{landscape}
\section*{D.3 Complete Variety Data (Ranks 1--101)}
All 101 grape varieties sorted by TasteRank. Columns show rank, variety name, type (Red / White), community assignment, TasteRank score, and the thirteen sensory profile dimensions.

\renewcommand{\arraystretch}{1.05}
\small
\begin{longtable}{@{}r l c c r *{13}{c}@{}}
\toprule
\# & Variety & Type & Com & TR & Col & Aro & Flo & Frt & Hrb & Spk & Acd & Tan & Bod & Alc & Fla & Fin & Cpx \\
\midrule
\endfirsthead
\toprule
\# & Variety & Type & Com & TR & Col & Aro & Flo & Frt & Hrb & Spk & Acd & Tan & Bod & Alc & Fla & Fin & Cpx \\
\midrule
\endhead
\bottomrule
\endlastfoot
1 & Sagrantino & Red & C0 & 0.3020 & 5 & 3 & 1 & 4 & 2 & 3 & 3 & 5 & 5 & 4 & 4 & 4 & 4 \\
2 & Nero d’Avola & Red & C0 & 0.2849 & 5 & 3 & 1 & 4 & 2 & 2 & 3 & 4 & 4 & 4 & 4 & 3 & 3 \\
3 & Lagrein & Red & C0 & 0.2748 & 5 & 3 & 1 & 3 & 2 & 2 & 3 & 4 & 4 & 3 & 3 & 3 & 3 \\
4 & Negroamaro & Red & C0 & 0.2380 & 4 & 3 & 1 & 4 & 2 & 2 & 3 & 3 & 4 & 4 & 3 & 3 & 3 \\
5 & Plavac Mali & Red & C0 & 0.2251 & 5 & 3 & 1 & 4 & 2 & 2 & 3 & 4 & 5 & 5 & 4 & 3 & 3 \\
6 & Montepulciano & Red & C0 & 0.2108 & 5 & 3 & 1 & 4 & 2 & 2 & 3 & 4 & 4 & 3 & 3 & 3 & 3 \\
7 & Petit Verdot & Red & C0 & 0.2093 & 5 & 3 & 2 & 3 & 2 & 3 & 3 & 5 & 5 & 4 & 4 & 4 & 4 \\
8 & Monastrell & Red & C0 & 0.2060 & 5 & 4 & 1 & 4 & 2 & 3 & 2 & 4 & 5 & 5 & 4 & 4 & 3 \\
9 & Petite Sirah & Red & C0 & 0.1768 & 5 & 3 & 1 & 4 & 1 & 3 & 3 & 5 & 5 & 4 & 4 & 4 & 3 \\
10 & Pinotage & Red & C0 & 0.1710 & 5 & 4 & 0 & 4 & 1 & 3 & 3 & 4 & 4 & 4 & 4 & 3 & 3 \\
11 & Malbec & Red & C0 & 0.1697 & 5 & 3 & 2 & 4 & 1 & 3 & 3 & 4 & 5 & 4 & 4 & 4 & 3 \\
12 & Tannat & Red & C0 & 0.1676 & 5 & 3 & 1 & 3 & 2 & 3 & 4 & 5 & 5 & 4 & 4 & 4 & 4 \\
13 & Garnacha Tintorera & Red & C0 & 0.1670 & 5 & 3 & 1 & 4 & 2 & 2 & 3 & 4 & 5 & 4 & 3 & 3 & 2 \\
14 & Dolcetto & Red & C0 & 0.1659 & 4 & 3 & 1 & 3 & 1 & 1 & 3 & 3 & 3 & 3 & 3 & 2 & 2 \\
15 & Mourvèdre & Red & C0 & 0.1565 & 5 & 4 & 1 & 3 & 3 & 3 & 3 & 4 & 5 & 4 & 4 & 4 & 4 \\
16 & Primitivo & Red & C0 & 0.1498 & 5 & 4 & 1 & 5 & 1 & 3 & 2 & 3 & 5 & 5 & 4 & 3 & 3 \\
17 & Zinfandel & Red & C0 & 0.1498 & 5 & 4 & 1 & 5 & 1 & 3 & 2 & 3 & 5 & 5 & 4 & 3 & 3 \\
18 & Merlot & Red & C0 & 0.1333 & 4 & 3 & 1 & 4 & 2 & 3 & 3 & 3 & 4 & 4 & 3 & 4 & 4 \\
19 & Carignan & Red & C0 & 0.1283 & 4 & 3 & 1 & 3 & 3 & 1 & 4 & 4 & 4 & 3 & 3 & 3 & 3 \\
20 & Limniona & Red & C2 & 0.1223 & 4 & 4 & 2 & 3 & 3 & 2 & 4 & 4 & 4 & 3 & 4 & 4 & 4 \\
21 & St. Laurent & Red & C2 & 0.1218 & 3 & 3 & 1 & 3 & 2 & 2 & 4 & 3 & 3 & 3 & 3 & 3 & 3 \\
22 & Bobal & Red & C0 & 0.1212 & 5 & 3 & 1 & 3 & 2 & 1 & 4 & 4 & 4 & 3 & 3 & 3 & 3 \\
23 & Grenache & Red & C0 & 0.1206 & 3 & 3 & 1 & 4 & 2 & 2 & 2 & 2 & 4 & 5 & 3 & 3 & 3 \\
24 & Syrah & Red & C0 & 0.1177 & 5 & 4 & 2 & 3 & 2 & 4 & 3 & 4 & 5 & 4 & 4 & 5 & 5 \\
25 & Tempranillo & Red & C0 & 0.1163 & 4 & 3 & 1 & 3 & 3 & 3 & 3 & 4 & 4 & 3 & 3 & 4 & 4 \\
26 & Mavrud & Red & C0 & 0.1149 & 5 & 3 & 1 & 3 & 3 & 2 & 4 & 5 & 4 & 4 & 4 & 4 & 4 \\
27 & Zweigelt & Red & C2 & 0.1134 & 3 & 3 & 1 & 3 & 1 & 1 & 3 & 2 & 3 & 3 & 3 & 3 & 2 \\
28 & Counoise & Red & C2 & 0.1121 & 3 & 3 & 2 & 3 & 2 & 1 & 4 & 3 & 3 & 3 & 3 & 3 & 3 \\
29 & Graciano & Red & C0 & 0.1027 & 5 & 4 & 2 & 3 & 3 & 3 & 4 & 5 & 4 & 3 & 4 & 4 & 4 \\
30 & Marsanne & White & C3 & 0.0978 & 3 & 3 & 1 & 3 & 1 & 2 & 2 & 0 & 4 & 4 & 3 & 3 & 3 \\
31 & Cabernet Sauvignon & Red & C0 & 0.0940 & 5 & 4 & 1 & 3 & 3 & 4 & 3 & 5 & 5 & 4 & 4 & 5 & 5 \\
32 & Tinta Roriz & Red & C0 & 0.0925 & 4 & 3 & 1 & 3 & 3 & 3 & 3 & 4 & 4 & 3 & 3 & 4 & 4 \\
33 & Blaufränkisch & Red & C2 & 0.0909 & 4 & 4 & 1 & 3 & 2 & 2 & 4 & 4 & 3 & 3 & 4 & 4 & 4 \\
34 & Fiano & White & C3 & 0.0855 & 2 & 4 & 2 & 3 & 1 & 1 & 4 & 0 & 3 & 3 & 4 & 4 & 4 \\
35 & Touriga Nacional & Red & C0 & 0.0849 & 5 & 5 & 3 & 3 & 2 & 3 & 3 & 5 & 5 & 4 & 5 & 5 & 5 \\
36 & Carménère & Red & C2 & 0.0839 & 4 & 4 & 1 & 3 & 4 & 2 & 3 & 3 & 4 & 3 & 4 & 3 & 3 \\
37 & Godello & White & C3 & 0.0826 & 2 & 3 & 1 & 3 & 1 & 1 & 4 & 0 & 3 & 3 & 3 & 3 & 4 \\
38 & Pinot Gris & White & C3 & 0.0803 & 3 & 3 & 1 & 3 & 0 & 1 & 3 & 0 & 3 & 3 & 3 & 3 & 3 \\
39 & Aglianico & Red & C0 & 0.0787 & 4 & 4 & 1 & 3 & 3 & 3 & 5 & 5 & 4 & 4 & 4 & 5 & 5 \\
40 & Baga & Red & C0 & 0.0787 & 4 & 3 & 1 & 2 & 3 & 2 & 5 & 5 & 3 & 3 & 3 & 4 & 4 \\
41 & Chardonnay & White & C3 & 0.0764 & 3 & 3 & 1 & 3 & 1 & 3 & 3 & 0 & 4 & 4 & 3 & 4 & 4 \\
42 & Sémillon & White & C3 & 0.0753 & 3 & 3 & 1 & 3 & 0 & 2 & 2 & 0 & 4 & 4 & 3 & 4 & 4 \\
43 & Encruzado & White & C3 & 0.0747 & 2 & 3 & 1 & 3 & 1 & 2 & 4 & 0 & 3 & 3 & 3 & 4 & 4 \\
44 & Greco & White & C4 & 0.0742 & 2 & 3 & 1 & 2 & 1 & 0 & 4 & 0 & 3 & 3 & 3 & 3 & 4 \\
45 & Roussanne & White & C3 & 0.0710 & 3 & 4 & 2 & 3 & 2 & 2 & 3 & 0 & 4 & 3 & 4 & 4 & 4 \\
46 & Friulano & White & C4 & 0.0688 & 2 & 3 & 1 & 2 & 1 & 0 & 3 & 0 & 3 & 3 & 3 & 3 & 3 \\
47 & Verdelho & White & C3 & 0.0683 & 2 & 3 & 1 & 3 & 0 & 1 & 3 & 0 & 3 & 3 & 3 & 3 & 3 \\
48 & Verdicchio & White & C4 & 0.0674 & 2 & 3 & 1 & 2 & 1 & 0 & 4 & 0 & 3 & 3 & 3 & 3 & 3 \\
49 & Cabernet Franc & Red & C0 & 0.0671 & 3 & 4 & 2 & 2 & 4 & 2 & 4 & 4 & 3 & 3 & 4 & 4 & 4 \\
50 & Verdejo & White & C4 & 0.0662 & 2 & 3 & 1 & 2 & 2 & 0 & 4 & 0 & 3 & 3 & 3 & 3 & 3 \\
51 & Barbera & Red & C2 & 0.0650 & 4 & 3 & 1 & 3 & 1 & 2 & 5 & 2 & 3 & 3 & 3 & 3 & 3 \\
52 & Vermentino & White & C4 & 0.0650 & 2 & 3 & 1 & 2 & 2 & 0 & 4 & 0 & 3 & 3 & 3 & 3 & 3 \\
53 & Corvina & Red & C2 & 0.0641 & 3 & 3 & 2 & 3 & 1 & 1 & 4 & 2 & 3 & 3 & 3 & 3 & 3 \\
54 & Sangiovese & Red & C2 & 0.0632 & 3 & 4 & 1 & 2 & 3 & 2 & 5 & 4 & 3 & 3 & 4 & 4 & 4 \\
55 & Mencía & Red & C2 & 0.0601 & 3 & 4 & 2 & 2 & 3 & 2 & 4 & 3 & 3 & 3 & 4 & 3 & 4 \\
56 & Grüner Veltliner & White & C4 & 0.0594 & 2 & 3 & 1 & 2 & 2 & 0 & 4 & 0 & 3 & 3 & 3 & 3 & 4 \\
57 & Pinot Blanc & White & C3 & 0.0574 & 2 & 2 & 1 & 2 & 0 & 1 & 3 & 0 & 3 & 3 & 2 & 3 & 2 \\
58 & Nebbiolo & Red & C2 & 0.0567 & 3 & 5 & 3 & 3 & 3 & 3 & 5 & 5 & 3 & 4 & 5 & 5 & 5 \\
59 & Arneis & White & C3 & 0.0560 & 2 & 3 & 2 & 3 & 0 & 0 & 3 & 0 & 3 & 3 & 3 & 3 & 3 \\
60 & Vidiano & White & C3 & 0.0550 & 2 & 3 & 2 & 3 & 1 & 1 & 3 & 0 & 3 & 3 & 3 & 3 & 3 \\
61 & Cinsault & Red & C2 & 0.0540 & 2 & 3 & 2 & 3 & 1 & 1 & 3 & 2 & 2 & 3 & 3 & 2 & 2 \\
62 & Touriga Franca & Red & C2 & 0.0524 & 4 & 4 & 3 & 3 & 1 & 2 & 3 & 3 & 4 & 3 & 4 & 4 & 4 \\
63 & Falanghina & White & C1 & 0.0515 & 2 & 3 & 2 & 2 & 0 & 0 & 4 & 0 & 3 & 3 & 3 & 3 & 3 \\
64 & Nerello Mascalese & Red & C2 & 0.0499 & 2 & 4 & 2 & 2 & 3 & 2 & 5 & 3 & 2 & 3 & 4 & 4 & 5 \\
65 & Garganega & White & C4 & 0.0497 & 2 & 2 & 1 & 2 & 1 & 0 & 3 & 0 & 3 & 3 & 2 & 3 & 3 \\
66 & Silvaner & White & C4 & 0.0487 & 2 & 2 & 1 & 2 & 1 & 0 & 3 & 0 & 3 & 3 & 2 & 3 & 3 \\
67 & Xinomavro & Red & C2 & 0.0481 & 3 & 4 & 1 & 2 & 3 & 2 & 5 & 5 & 3 & 3 & 4 & 5 & 5 \\
68 & Clairette & White & C4 & 0.0437 & 2 & 2 & 1 & 2 & 1 & 0 & 3 & 0 & 3 & 3 & 2 & 2 & 2 \\
69 & Assyrtiko & White & C1 & 0.0409 & 2 & 4 & 1 & 2 & 1 & 0 & 5 & 0 & 3 & 3 & 4 & 4 & 5 \\
70 & Viognier & White & C3 & 0.0400 & 3 & 5 & 4 & 4 & 0 & 2 & 2 & 0 & 4 & 4 & 4 & 3 & 4 \\
71 & Pinot Noir & Red & C2 & 0.0399 & 2 & 4 & 2 & 2 & 2 & 2 & 4 & 2 & 2 & 3 & 4 & 4 & 5 \\
72 & Chenin Blanc & White & C4 & 0.0391 & 2 & 3 & 2 & 3 & 1 & 1 & 5 & 0 & 3 & 3 & 3 & 4 & 4 \\
73 & Albariño & White & C1 & 0.0387 & 1 & 4 & 2 & 2 & 1 & 0 & 5 & 0 & 2 & 3 & 3 & 3 & 4 \\
74 & Gewürztraminer & White & C5 & 0.0361 & 3 & 5 & 4 & 4 & 1 & 2 & 2 & 0 & 4 & 4 & 5 & 3 & 4 \\
75 & Furmint & White & C1 & 0.0361 & 2 & 3 & 1 & 2 & 1 & 1 & 5 & 0 & 3 & 3 & 3 & 4 & 5 \\
76 & Petit Manseng & White & C5 & 0.0356 & 2 & 4 & 2 & 4 & 1 & 1 & 5 & 0 & 3 & 3 & 4 & 4 & 4 \\
77 & Riesling & White & C1 & 0.0353 & 1 & 4 & 3 & 2 & 0 & 0 & 5 & 0 & 2 & 2 & 4 & 4 & 5 \\
78 & Trousseau & Red & C2 & 0.0352 & 2 & 3 & 1 & 2 & 2 & 1 & 4 & 2 & 2 & 2 & 3 & 3 & 3 \\
79 & Kadarka & Red & C1 & 0.0335 & 2 & 3 & 2 & 2 & 2 & 1 & 4 & 2 & 2 & 2 & 3 & 2 & 3 \\
80 & Loureiro & White & C1 & 0.0324 & 1 & 4 & 3 & 2 & 0 & 0 & 4 & 0 & 2 & 2 & 3 & 2 & 3 \\
81 & Ribolla Gialla & White & C1 & 0.0318 & 1 & 3 & 1 & 2 & 1 & 0 & 4 & 0 & 2 & 2 & 3 & 2 & 3 \\
82 & Malvasia & White & C5 & 0.0315 & 2 & 4 & 3 & 3 & 1 & 0 & 3 & 0 & 3 & 3 & 3 & 3 & 3 \\
83 & Gamay & Red & C2 & 0.0310 & 2 & 4 & 2 & 2 & 1 & 1 & 4 & 1 & 2 & 2 & 3 & 2 & 3 \\
84 & Muscat Blanc & White & C5 & 0.0302 & 2 & 5 & 5 & 3 & 0 & 0 & 3 & 0 & 2 & 2 & 4 & 2 & 3 \\
85 & Malagousia & White & C5 & 0.0297 & 2 & 4 & 4 & 3 & 1 & 0 & 3 & 0 & 3 & 3 & 4 & 3 & 4 \\
86 & Torrontés & White & C5 & 0.0296 & 2 & 5 & 4 & 3 & 0 & 0 & 3 & 0 & 3 & 3 & 4 & 2 & 3 \\
87 & Sauvignon Blanc & White & C1 & 0.0290 & 1 & 4 & 1 & 2 & 4 & 0 & 5 & 0 & 2 & 2 & 4 & 3 & 3 \\
88 & Cortese & White & C1 & 0.0256 & 1 & 2 & 1 & 2 & 0 & 0 & 4 & 0 & 2 & 2 & 2 & 2 & 2 \\
89 & Arinto & White & C1 & 0.0255 & 1 & 3 & 1 & 2 & 0 & 0 & 5 & 0 & 2 & 2 & 3 & 3 & 3 \\
90 & Welschriesling & White & C1 & 0.0234 & 1 & 3 & 2 & 2 & 0 & 0 & 4 & 0 & 2 & 2 & 2 & 2 & 2 \\
91 & Müller-Thurgau & White & C1 & 0.0226 & 1 & 3 & 2 & 2 & 0 & 0 & 3 & 0 & 2 & 2 & 2 & 2 & 2 \\
92 & Schiava & Red & C1 & 0.0219 & 1 & 3 & 2 & 2 & 1 & 0 & 4 & 1 & 1 & 2 & 2 & 2 & 2 \\
93 & Poulsard & Red & C1 & 0.0219 & 1 & 3 & 2 & 2 & 1 & 0 & 4 & 1 & 1 & 2 & 2 & 2 & 2 \\
94 & Colombard & White & C1 & 0.0214 & 1 & 2 & 0 & 2 & 1 & 0 & 4 & 0 & 2 & 2 & 2 & 2 & 2 \\
95 & Chasselas & White & C5 & 0.0209 & 1 & 2 & 0 & 2 & 0 & 0 & 3 & 0 & 2 & 2 & 2 & 2 & 2 \\
96 & Trebbiano & White & C5 & 0.0196 & 1 & 2 & 0 & 2 & 0 & 0 & 3 & 0 & 2 & 2 & 2 & 2 & 2 \\
97 & Frappato & Red & C1 & 0.0195 & 2 & 4 & 3 & 2 & 1 & 0 & 4 & 1 & 2 & 2 & 3 & 2 & 3 \\
98 & Aligoté & White & C1 & 0.0188 & 1 & 2 & 0 & 2 & 0 & 0 & 4 & 0 & 2 & 2 & 2 & 2 & 2 \\
99 & Melon de Bourgogne & White & C1 & 0.0174 & 1 & 2 & 0 & 1 & 0 & 0 & 4 & 0 & 2 & 2 & 2 & 2 & 2 \\
100 & Mauzac & White & C1 & 0.0167 & 1 & 2 & 1 & 2 & 0 & 0 & 4 & 0 & 2 & 2 & 2 & 2 & 2 \\
101 & Ugni Blanc & White & C5 & 0.0101 & 1 & 1 & 0 & 1 & 0 & 0 & 5 & 0 & 1 & 2 & 1 & 1 & 1 \\

\end{longtable}
\end{landscape}

\section*{References}

\begin{itemize}[leftmargin=2em,itemsep=0.2em]
\item d'Agata, I. (2014). \emph{Native Wine Grapes of Italy}. University of California Press.
\item Galet, P. (2000). \emph{Dictionnaire encyclop\'edique des c\'epages}. Hachette, Paris.
\item Peynaud, \'E. (1980). \emph{Le Go\^ut du Vin}. Dunod, Paris.
\item Robinson, J., Harding, J. \& Vouillamoz, J. (2012). \emph{Wine Grapes: A Complete Guide to 1{,}368 Vine Varieties}. Ecco / HarperCollins.
\end{itemize}

\end{document}
